% Generated draft: requires mathematical and source review.
\documentclass[11pt]{article}
\usepackage[margin=25mm]{geometry}
\usepackage{fontspec}
\setmainfont{Latin Modern Roman}
\usepackage{amsmath,amssymb,mathtools,unicode-math}
\setmathfont{Latin Modern Math}
\usepackage{enumitem,needspace,xurl,hyperref}
\setlength{\parindent}{0pt}
\setlength{\parskip}{5pt}
\setlength{\emergencystretch}{3em}
\newcommand{\N}{\mathbb{N}}
\newcommand{\Z}{\mathbb{Z}}
\newcommand{\Q}{\mathbb{Q}}
\newcommand{\R}{\mathbb{R}}
\newcommand{\C}{\mathbb{C}}
\newcommand{\F}{\mathbb{F}}
\newcommand{\A}{\mathbb{A}}
\newcommand{\PP}{\mathbb{P}}
\newcommand{\E}{\mathbb{E}}
\newcommand{\eps}{\varepsilon}
\newcommand{\dd}{\,\mathrm{d}}
\begin{document}
\section*{UnsolvedMath ID: 30003261}
\textbf{Source problem code:} catalogue entry
\subsection{Definitions and mathematical statement}
The supplied record is a composite extraction marked \emph{unrecoverable}; it does not determine a single fully specified mathematical problem. Its two components are retained below, with unresolved choices stated explicitly.

Let \(k\) be an algebraically closed complete valued field. The record denotes by \(A_k\) a rigid entire-function ring and by \(M_k\) a rigid meromorphic field, with the coordinate \(t\) distinguished in each. It does not supply definitions of these analytic objects or conventions concerning the valuation. These symbols are therefore retained as source notation, pending recovery of their precise definitions.

For a ring or field \(S\) with a distinguished element \(t_S\), let \(\mathcal L_t=\{0,1,+,-,\cdot,t\}\), where the constant symbol \(t\) is interpreted as \(t_S\). Define
\[
\operatorname{Th}_{\exists,\mathcal L_t}(S,t_S)
 =\left\{\exists x_1\cdots\exists x_n\,\varphi(x_1,\ldots,x_n):
 \begin{array}{l}
 n\geq 0,\ \varphi\text{ is quantifier-free in }\mathcal L_t,\\
 (S,t_S)\models\exists x_1\cdots\exists x_n\,\varphi
 \end{array}\right\},
\]
and let \(\operatorname{Th}_{\mathcal L_t}(S,t_S)\) be the set of all first-order \(\mathcal L_t\)-sentences true in \((S,t_S)\). A theory is decidable if an algorithm, on every input sentence of its language, terminates and correctly determines membership in the theory. The language \(\mathcal L_t\) makes explicit the minimal ring-language interpretation of ``with \(t\) named''; the record does not establish whether additional constants, predicates, or sorts were intended.

\textbf{Analytic component.} Determine the decidability of
\[
\operatorname{Th}_{\exists,\mathcal L_t}(A_k,t)
\]
in each of the following cases:
\begin{enumerate}
\item \(k\) is archimedean;
\item \(k\) is non-archimedean and \(\operatorname{char}(k)=0\);
\item \(k\) is non-archimedean and \(\operatorname{char}(k)>0\).
\end{enumerate}
The record also asks for decidability of ``the theory'' of \((M_k,t)\) in the same three cases. It does not resolve whether this means
\(\operatorname{Th}_{\mathcal L_t}(M_k,t)\) or only
\(\operatorname{Th}_{\exists,\mathcal L_t}(M_k,t)\); these are distinct targets and neither is silently selected here. The question is retained for an arbitrary \(k\) satisfying the stated hypotheses. Whether algorithms must be uniform as \(k\) varies, and how any effective presentation of \(k\) enters the question, are unspecified.

\textbf{Separate arithmetic component.} The record asks which \(p\)-adic analogues of undecidability results for rings of totally real algebraic integers hold, and mentions as an example ``the ring of integers of \(\Q_p^t\) where \(p\) is totally split.'' Here \(p\) is a prime. The field denoted by \(\Q_p^t\), the base extension in which total splitting is meant, the meaning of its ring of integers, the logical language, and the particular analogue sought are not defined in the supplied record. Thus this component is a retained research question, not a uniquely specified decision problem; no identification of \(\Q_p^t\) with a particular field is asserted.

The supplied verification note says that the official problem-session list separates these two components and that no clean statement was promoted. The supplied dated assessment labels the composite \emph{partially solved}, reports known undecidability in several listed cases while leaving pure ring theories and other cases open, and qualifies the cited 2023 \(p\)-adic result as using a richer valued-field language that does not settle the pure-ring questions. It does not identify all affected cases. These are supplied curation claims, not independently verified mathematical conclusions.
\subsection{Short English statement}
This composite record asks about decidability for entire-function rings and meromorphic fields over algebraically closed complete valued fields, and separately asks for \(p\)-adic analogues of undecidability for totally real algebraic integers. Missing definitions and language choices prevent a unique precise formulation. The supplied \emph{partially solved} label requires independent review.
\subsection{Sources}
\begin{itemize}
\item Local UnsolvedMath registry: \texttt{data/unsolvedmath/problems.json}, permanent ID \(30003261\), source problem code \texttt{catalogue entry}, title ``Composite of Unrelated Decidability Questions.'' The supplied record gives proposed year 2016 and names no proposer. Its statement status is \texttt{unrecoverable}; its supplied verification note says the extraction joins separate questions and that no clean statement was promoted.
\item Supplied report citation: \emph{Definability and Decidability Problems in Number Theory} (2017), Oberwolfach Reports, pp.\ 2793--2866, DOI \texttt{10.4171/owr/2016/49}. Source URL and literature-source label ``OWR 49 report'': \url{https://doi.org/10.4171/owr/2016/49}. The background attributes the material to workshop open-problem discussions documented in Oberwolfach Reports.
\item Supplied literature source ``Kartas 2023'': \url{https://doi.org/10.1016/j.apal.2022.103203}. The supplied assessment qualifies the relevant \(p\)-adic result as using a richer valued-field language and not settling pure-ring questions.
\item Supplied literature source ``Fehm 2015'': \url{https://arxiv.org/abs/1501.04522}.
\item Literature assessment checked at \texttt{2026-08-21}: the supplied status is \emph{partially solved}. The assessment reports undecidability in several listed cases and unresolved pure ring theories and other cases, without a complete case-by-case allocation. This is supplied dated curation derived from the OWR material, not independently verified here; its time-sensitive status and cited claims require independent verification.
\end{itemize}
\end{document}
